\section{Discussion and limitations}\label{sec:discussion}

\subsection{What the frontier does and does not settle}

A certified frontier answers a budgeting question, not a doctrine question. It
says which revision calendars are compatible with a tolerated loss under the
frozen model, and it carries the residual of the computation that produced it.
It does not say that a particular calendar is optimal in the real world, and
it does not identify the mechanism by which revisions pay. Two readings of the
same frontier should be resisted.

First, monotonicity in $K$ is a correctness property, not a finding: it follows
from policy inclusion (Proposition~\ref{prop:inclusion}) and would hold for any
model in this class. The informative content is the \emph{shape}: where the
frontier saturates, and which epochs the maximizing calendar selects.

Second, a budget that suffices is not a budget that is necessary. The
deletion certificate bounds the loss of a specified coarse calendar from
above; the interval certificate bounds a budget by a shortest path over
interval losses. Neither claims minimality for a true performance target, and
neither should be read as an optimality certificate for the calendar itself.

\subsection{Where the certificates can be loose}

Both certificates can be conservative, and we say where. The interval bound
takes a maximum over all histories at a boundary
(\eqref{eq:intervalcost}), so one unfavourable history sets the cost of an
interval for every path through it. The deletion certificate multiplies source
transition kernels along hypothetical completions of the opponent's future,
which is exact for the constructed policy but says nothing about the best
coarse policy. When either bound is far above the realised loss, the honest
report is that the certificate is weak in that cell, not that the budget is
tight. The certificates are also asymmetric in the opponent class: the
interval construction needs public boundaries, which an unexecuted committed
opponent suffix destroys, so it applies to the fully flexible opponent only.

\subsection{Model scope}

The numerical instantiation is a leader--follower line-ahead formation with
fixed arc-length spacing, a frozen directional kernel calibrated against an
auditable rules engine, and a fixed horizon with unit discount. Four
restrictions bound the conclusions. (i)~Formations are rigid in shape: no
formation breaking, no independent vessel autonomy, and no command delay.
(ii)~Information is perfect within the model's observation structure; the
budget restricts \emph{when} a player may act, not what it can see, so this
paper is not about costly observation. (iii)~The payoff is a reduced
interaction kernel, not a physical ballistic or hydrodynamic model; it is a
surrogate whose fidelity was measured separately and whose transfer to the
rules engine is partial. (iv)~The design is a fixed computational benchmark:
conditions were declared before their outputs, the solver is deterministic,
and no statistical sampling claim is attached to its output. Where a value is
unresolved, we report the interval and the termination reason rather than a
point estimate.

\subsection{Falsified predecessor and its role}

An earlier iteration of this project asked whether the side with the range
advantage values revision rights differently from the side without it. Tested
at this fidelity with exact anchors and frozen held-out configurations, that
claim failed: most configurations pointed the other way, and neither the sign
nor the ranking was stable. We keep the negative result in the record, do not
re-tune the design to revive it, and do not restate it under a new name. The
counterexample of Proposition~\ref{prop:counterexample} explains why the
underlying intuition --- that more revision opportunities must have
diminishing value --- cannot be assumed: revision opportunities can be
complementary, so their value is a property of the instance.

\subsection{What would move this to a stronger claim}

Three things would strengthen the paper beyond its present scope, and we list
them as open rather than claim them. A certificate that is tight rather than
merely valid, for instance by computing interval losses on the realised
history class instead of the worst history. A budgeted equilibrium notion in
which both players choose calendars, since the present design fixes one
calendar as public and optimises the other. And a formation model with
deformation, which would test whether the frontier shape survives the
representation change that the ablation probes only at $T=4$.
